% ======================================================================
% Ontologie der Kontinua -- Core 1.3.3
% Cross-K-Modul: crossk_k0_k1.tex
% Cross-Level-Relationen zwischen K0 und K1
% ======================================================================

\subsubsection{K0–K1-Querverbindung}
\label{sec:crossk-k0-k1}

Der Übergang von $K_0$ zu $K_1$ repräsentiert die erste und fundamentalste
Cross-Level-Struktur in der gesamten Hierarchie der Kontinua. Er beschreibt die
Geburt unterscheidbarer Zustände, das Auftreten einer expliziten Achse und die
Bildung eines nicht-trivialen Zustandsraums $\Omega(K_1)$ aus dem minimalen
prä-ontologischen Substrat von $K_0$.

Dieses Kapitel fasst die strukturellen Abhängigkeiten, vererbten Komponenten,
Übergangsoperatoren und Schwellenbedingungen für das Auftreten von $K_1$ zusammen.

% ----------------------------------------------------------------------
\subsubsection{1. Beteiligte Stufen und ihre Rollen}

\paragraph{K0.}
$K_0$ ist durch das \emph{Axiom der Differenz und Verbundenheit} definiert.
Es enthält:
\begin{itemize}
  \item einen minimalen Raum potenzieller Zustände ohne explizite Achsen;
  \item ein strukturelles Differenzmaß $\Delta(s_1,s_2)$;
  \item die Existenzschwelle $\Theta_0$ zur Sicherung der Trivialitätsfreiheit:
        $0 < \Delta < \varepsilon$ für ein $\varepsilon>0$;
  \item keine Zeit, keine Flüsse $J$, keine Zyklen $C$, keine Dynamik.
\end{itemize}

\paragraph{K1.}
$K_1$ ist das erste Kontinuum mit:
\begin{itemize}
  \item einer expliziten Achse $A_1$, die eine stabile Klasse von Unterschieden
        erfasst;
  \item einem definierten Zustandsraum $\Omega(K_1)$ mit Rand
        $\partial\Omega(K_1)$;
  \item einer primitiven Energie-/Potentialstruktur $P_1$;
  \item minimalen Flüssen $J_1$ und einem trivialen Zyklus
        $C_{\mathrm{triv}}$;
  \item expliziten Schwellen $\Theta_1$ (Existenz, Gradienten, Stabilität).
\end{itemize}

Die Rolle des Übergangs K0→K1 ist die \textbf{Erzeugung von Struktur}:
die Ausbildung der ersten darstellbaren Dimension im Kontinuum.

% ----------------------------------------------------------------------
\subsubsection{2. Gemeinsame und vererbte Achsen und Schwellen}

Obwohl $K_0$ keine Achsen besitzt, werden mehrere seiner strukturellen
Eigenschaften in transformierter Form von $K_1$ vererbt:

\paragraph{Differenz \texorpdfstring{$\Delta$}{\Delta}.}
Die primitive Differenzstruktur von $K_0$ wird zur metrischen Komponente der
Achse $A_1$ in $K_1$:
\[
A_1 \sim f(\Delta),
\]
womit die Existenz von Unterscheidbarkeit in $K_0$ die Entstehung einer
expliziten Koordinate in $K_1$ ermöglicht.

\paragraph{Schwellen.}
Die Existenzschwelle $\Theta_0$ wird zu einer \emph{Untergrenze} in $K_1$:
Wenn $\Theta_0$ verletzt wird (d.\,h. $\Delta=0$), kann keine Achse erzeugt werden
und $\Omega(K_1)=\varnothing$.

Entsprechend führt $\Theta_1$ neue Stabilitätsbedingungen für Gradienten,
Grenzen und zulässige Flüsse ein. Diese Schwellen gibt es in $K_0$ nicht und sie
entstehen erst nach dem Auftreten der ersten Achse.

% ----------------------------------------------------------------------
\subsubsection{3. Cross-Level-Flüsse, Zyklen und Spannungen}

\paragraph{Flüsse.}
In $K_0$ gibt es keine intrinsischen Flüsse, daher hat der Operator $J_1$ von
$K_1$ keinen Vorläufer. Stattdessen entsteht er aus der neu verfügbaren
Struktur:
\[
J_1 : A_1 \rightarrow \Omega(K_1)
\]
und stellt Bewegungen entlang oder quer der neuen Achse dar.

\paragraph{Zyklen.}
$K_0$ kann keine Zyklen tragen, da es keine Zeit und keine Flüsse besitzt.
Der triviale Zyklus von $K_1$:
\[
C_{\mathrm{triv}} : s \mapsto s,
\]
ist die erste mögliche zyklische Struktur in der Hierarchie.

\paragraph{Spannung.}
Die Cross-Level-strukturelle Spannung zwischen $K_0$ und $K_1$ ist definiert durch:
\[
T_{0,1} = g(\Delta - \Theta_0,\ P_1 - \Theta_1).
\]
Wenn $T_{0,1}$ seine dimensionsbezogene Schwelle $\Theta_{1,\mathrm{dim}}$
überschreitet, wird die entstehende Achse instabil oder kollabiert.

% ----------------------------------------------------------------------
\subsubsection{4. Geburts- und Todesbedingungen zwischen den Stufen}

\paragraph{Geburt von \texorpdfstring{$K_1$}{K_1}.}
$K_1$ tritt auf, wenn:
\[
\Delta > \Theta_0
\quad \text{und} \quad
T_0 > \Theta_{0,\mathrm{dim}},
\]
was das Erscheinen einer neuen Achse $A_1$ erlaubt:
\[
A_1 \in M_1 \setminus A(K_0).
\]

Dies ist die Wirkung des Operators $\Psi_{0\rightarrow 1}$ im Core:
Er erweitert die minimale Struktur von $K_0$ zu einem repräsentationalen
Kontinuum $K_1$.

\paragraph{Todesausbreitung.}
Wenn $\Theta_0$ versagt, kann $K_1$ nicht existieren:
\[
\Theta_0^{\mathrm{death}}
\Rightarrow \Omega(K_1)=\varnothing.
\]

Wenn $\Theta_1$ versagt, kollabiert das Kontinuum $K_1$, während $K_0$ intakt
bleibt, da $K_0$ keine strukturellen Abhängigkeiten von $K_1$ hat.

% ----------------------------------------------------------------------
\subsubsection{5. Konzeptionelle Beispiele}

\paragraph{Entstehung der ersten Achse aus minimaler Differenz.}
Jede Situation, in der eine rohe Unterscheidung (z.\,B. ``Zustand A unterscheidet
sich von Zustand B'') stabilisiert und parametrisiert werden kann, entspricht
einer Realisierung des Übergangs K0→K1. Die Achse $A_1$ verkörpert diese
stabilisierte Unterscheidung.

\paragraph{Randbildung.}
Sobald die Achse vorhanden ist, werden Grenzen sinnvoll:
\[
\partial\Omega(K_1) = \{ s \in \Omega(K_1) \mid \nabla A_1 \text{ erreicht eine Schwelle} \}.
\]

\paragraph{Erste Flüsse.}
Mit einer Achse und Grenzbildung ergeben sich gradientenähnliche Flüsse:
eine Unmöglichkeit in $K_0$, aber eine natürliche Folge der stabilisierten
Unterscheidung in $K_1$.

Diese Beispiele veranschaulichen die globale Bedeutung des K0→K1-Übergangs:
\[
\textbf{Es ist die Geburt aller darstellbaren Struktur in der
Kontinuums-Hierarchie.}
\]
